\chapter{Bitcoin, USDT, Collection Partners, and Payout Settlement}
\label{ch:btc-usdt-settlement}
\chapterline

At the end of Chapter~\ref{ch:fx-corridors}, Savanna Pay faced a costly bottleneck. Paying Mary in Nairobi took ten seconds only because Savanna Pay kept a large pile of Kenyan shillings sitting idle inside its Nairobi tub. Refilling that tub through traditional SWIFT bank wires took three to four days and paused completely every Friday afternoon until Monday morning.

Those slow wires forced Savanna Pay to lock up hundreds of thousands of dollars in dead working capital and exposed its thin $1.8\%$ fee margin to weekend currency swings.

One afternoon in Nairobi, Tunde asks a practical engineering question. ``Why are we waiting four days for correspondent banks when a stablecoin transfer settles globally in thirty minutes, even at midnight on Saturday?''

Today, many cross-border payment companies settle their wholesale treasury leg over stablecoin rails in under half an hour. Even better, neither the sender in New York or London nor the grandmother in Nairobi ever touches cryptocurrency. Let us build how this works from zero, compare Bitcoin and dollar stablecoins honestly, and trace how every book settles to the exact cent.

\section{Bitcoin vs Dollar Stablecoins From First Principles}

If you have never worked with a blockchain, strip away the hype and picture a shared, append-only ledger replicated across thousands of computers worldwide. Unlike a commercial bank ledger that closes at 17:00 on Friday, a public blockchain processes transfers 24 hours a day, 365 days a year.

Two very different digital assets run on public blockchains as cross-border payment rails: \textbf{Bitcoin (\texttt{BTC})} and \textbf{dollar stablecoins} such as \textbf{USDT (Tether)} and \textbf{USDC (Circle)}.

\begin{historybox}
In January 2009, Satoshi Nakamoto launched \textbf{Bitcoin}, the first digital asset that could move globally from person to person without a central bank or clearinghouse. Anyone holding a cryptographic private key could send value across borders on a Sunday night. For a remittance business earning half a percent per transfer, however, Bitcoin had a major catch: its market price in US dollars routinely swung $2\%$ to $5\%$ in a single afternoon.

To keep 24/7 blockchain speed while removing price swings, companies created \textbf{stablecoins} (blockchain tokens designed to hold a steady $1.00\text{ USD}$ price by backing every issued token with cash and short-term US government Treasury bills). Tether launched \textbf{USDT} in 2014, and Circle launched \textbf{USDC} in 2018. Running on fast networks such as Tron, Solana, and Ethereum Layer-2 chains, dollar stablecoins quickly became the wholesale cargo rail of emerging-market treasury desks.
\end{historybox}

\subsection{Comparing BTC and USDT as Treasury Plumbing}

Before wiring company money across a blockchain, an engineer and an accountant must check four practical properties (Table~\ref{tab:btc-vs-stablecoin}):
\begin{enumerate}
  \item \textbf{Price volatility during transit:} Suppose Savanna Pay earns $45\text{ bps}$ ($0.45\%$) net profit on a $50{,}000\text{ USD}$ batch ($225\text{ USD}$). If Savanna Pay buys $50{,}000\text{ USD}$ of unhedged Bitcoin, waits $40\text{ minutes}$ to convert it in Nairobi, and Bitcoin dips by a routine $1.5\%$, Savanna Pay loses $750\text{ USD}$ -- wiping out the entire $225\text{ USD}$ profit. Therefore, Bitcoin only works for fiat-to-fiat corridors when swapped in seconds over the \textbf{Lightning Network} (a fast payment-channel layer built on top of Bitcoin) or hedged simultaneously. By contrast, \textbf{USDT} and \textbf{USDC} trade tightly near $1.00\text{ USD}$ (typically within $0.999$ to $1.001\text{ USD}$), so holding them for thirty minutes or a weekend carries almost zero exchange-rate swing against the dollar.
  \item \textbf{Blockchain finality:} In crypto, \textbf{finality} (or \textbf{confirmations}) means waiting until enough new blocks are added that a transfer can never be reversed. Bitcoin's base layer mines one block roughly every $10\text{ minutes}$, so desks wait $30\text{ to }60\text{ minutes}$ for three to six confirmations. Fast stablecoin chains confirm in $5\text{ to }60\text{ seconds}$.
  \item \textbf{Network gas fees:} Every blockchain charges a small processing fee called a \textbf{network gas fee} to include your transfer in a block. Crucially, sending \texttt{USDT} does \emph{not} charge the fee in \texttt{USDT}; you must pay gas in the host blockchain's native coin (\texttt{TRX} on Tron, \texttt{ETH} on Ethereum or Base, \texttt{SOL} on Solana).
  \item \textbf{Issuer and freeze risk:} Nobody controls Bitcoin's base protocol, so no company can freeze an unencumbered Bitcoin wallet. By contrast, Tether and Circle hold off-chain bank and Treasury reserves and maintain smart-contract admin keys that can \textbf{freeze} \texttt{USDT} or \texttt{USDC} at addresses flagged by law enforcement or sanctions regulators.
\end{enumerate}

\begin{table}[htbp]
\centering
\caption{Comparing Bitcoin (\texttt{BTC}) and dollar stablecoins (\texttt{USDT} / \texttt{USDC}) for cross-border treasury settlement.}
\label{tab:btc-vs-stablecoin}
\small
\setlength{\tabcolsep}{4.5pt}
\begin{tabular}{p{2.5cm}p{4.9cm}p{5.0cm}}
\toprule
\textbf{Property} & \textbf{Bitcoin (\texttt{BTC})} & \textbf{Dollar Stablecoins (\texttt{USDT} / \texttt{USDC})} \\
\midrule
\textbf{Dollar Price Swing} & Swings $1\%$ to $4\%$ daily; risky to hold unhedged even for an hour. & Pegged near \$1.00 (usually within $\pm 0.05\%$ to $0.15\%$). \\
\addlinespace
\textbf{Confirmation Speed} & $30\text{--}60\text{ min}$ on base layer; under $2\text{ seconds}$ on Lightning Network. & $5\text{--}60\text{ seconds}$ on Tron, Solana, Base, or Arbitrum; $2\text{--}12\text{ min}$ on Ethereum L1. \\
\addlinespace
\textbf{Network Fee Coin} & Paid directly in \texttt{BTC} (tiny on Lightning; variable on base layer). & Paid in the host chain's native coin (\texttt{TRX}, \texttt{ETH}, or \texttt{SOL}), not in \texttt{USDT}. \\
\addlinespace
\textbf{Counterparty Risk} & Zero central issuer; cannot be frozen on the base protocol. & Relies on issuer's T-bill reserves; issuer can freeze sanctioned addresses. \\
\bottomrule
\end{tabular}
\end{table}

\begin{trapbox}
\textbf{The Empty Gas Tank Outage on Saturday Night.} Here is a classic stablecoin bug. Savanna Pay's automated treasury wallet holds $150{,}000\text{ USDT}$ on Saturday at midnight and tries to refill Nairobi. Every transfer fails with an obscure node error because the wallet ran out of $2.00\text{ USD}$ worth of native \texttt{TRX} or \texttt{ETH} needed to pay network gas. Always track a dedicated \texttt{Crypto Gas Inventory} account on your ledger and alert on-call engineers whenever native gas drops below a week of transfers.
\end{trapbox}

\section{The Real Two-Provider Architecture}

No single bank lets Savanna Pay charge debit cards in New York and pay out instant M-Pesa in Nairobi, MTN MoMo in Kampala, and Airtel Money in Accra. Instead, Savanna Pay connects three specialized partners into one clean pipeline (Figure~\ref{fig:four-ledger-settlement}):

\begin{enumerate}
  \item \textbf{Collection Partner (Origin PSP in the US or UK).} A \textbf{PSP} (\textbf{Payment Service Provider}: a regulated processor such as Stripe, Checkout, or an ACH bank gateway) collects card and bank payments from senders. The PSP deducts its fee and holds funds in a clearing bucket for \textbf{T+2} days before wiring cash to Savanna Pay's US bank. On card traffic, the PSP also holds a \textbf{rolling reserve} (a $5\%$ cash slice kept for $90$ days to cover \textbf{chargebacks}, when a cardholder disputes a charge and their bank pulls the money back).
  \item \textbf{Treasury On-Ramp (Regulated US/UK OTC Desk).} An \textbf{OTC desk} (\textbf{over-the-counter desk}: a licensed institutional dealer trading large blocks at a locked quote) swaps Savanna Pay's USD wire into \texttt{USDT}. Converting bank cash into crypto tokens is called \textbf{on-ramping}, and converting tokens back into bank cash is called \textbf{off-ramping}.
  \item \textbf{Payout Orchestrator (HoneyCoin).} A \textbf{payout orchestrator} is a regional payment hub wired directly into local mobile-money networks and banks across many countries. Platforms like \textbf{HoneyCoin} connect one API to M-Pesa in Kenya and MTN MoMo and Airtel Money across Uganda, Ghana, Tanzania, and Nigeria. Savanna Pay sends \texttt{USDT} on-chain to its treasury wallet at HoneyCoin. When a transfer passes compliance, Savanna Pay calls HoneyCoin's payout API, which debits Savanna Pay's float at an agreed shilling rate plus rail fee and delivers instant \texttt{KES} to the recipient's M-Pesa phone.
\end{enumerate}

\begin{happeningbox}
Notice who touches crypto and who does not. A sender in New York pays $200\text{ USD}$ with an ordinary Visa debit card or bank transfer. Ten seconds later, Mary in Nairobi receives $25{,}532\text{ KES}$ in her ordinary M-Pesa wallet. Neither customer ever downloads a crypto wallet or reads a blockchain address. \texttt{USDT} sits strictly under the hood as a thirty-minute wholesale cargo ship refilling Savanna Pay's HoneyCoin float pool.
\end{happeningbox}

\begin{figure}[htbp]
\centering
\begin{tikzpicture}[x=1cm,y=1cm,>=Latex,font=\scriptsize]
  % Lane backgrounds
  \draw[draw=sourceblue!50, fill=sourceblue!4, rounded corners=4pt] (0,2.55) rectangle (12.4,4.75);
  \node[font=\bfseries\scriptsize, text=sourceblue!90!black, anchor=west] at (0.2,4.48) {1. Instant Customer Path ($t = 0\text{ to }10\text{ seconds}$)};

  \draw[draw=quotientorange!60, fill=quotientorange!5, rounded corners=4pt] (0,0.0) rectangle (12.4,2.25);
  \node[font=\bfseries\scriptsize, text=quotientorange!85!black, anchor=west] at (0.2,1.98) {2. Wholesale Book Settlement Path (\texttt{T+2} Card Clear + $30\text{ min}$ USDT Refill)};

  % Top lane nodes
  \node[draw=gray!80!black, fill=white, rounded corners=3pt, minimum width=2.0cm, minimum height=0.95cm, align=center] (cust) at (1.25,3.45) {\textbf{US Sender}\\Marcus Pays\\\$200 / \$50k Batch};
  \node[draw=sourceblue!85!black, thick, fill=white, rounded corners=3pt, minimum width=2.3cm, minimum height=0.95cm, align=center] (pspa) at (4.3,3.45) {\textbf{Collection PSP}\\Card / Bank Rail\\($80\text{ bps}$ Fee)};
  \node[draw=structurepurple!90!black, thick, fill=white, rounded corners=3pt, minimum width=2.3cm, minimum height=0.95cm, align=center] (hcapi) at (7.85,3.45) {\textbf{HoneyCoin Float}\\Pre-Funded KES\\($40\text{ bps}$ Rail Fee)};
  \node[draw=targetgreen!85!black, thick, fill=white, rounded corners=3pt, minimum width=2.1cm, minimum height=0.95cm, align=center] (mpesa) at (11.1,3.45) {\textbf{KE Recipient}\\Mary's M-Pesa\\Instant KES};

  \draw[->, thick, sourceblue] (cust) -- (pspa);
  \draw[->, thick, dashed, structurepurple] (pspa) -- node[above, font=\tiny\bfseries] {Webhook $\to$ API} (hcapi);
  \draw[->, thick, targetgreen!80!black] (hcapi) -- (mpesa);

  % Bottom lane nodes
  \node[draw=sourceblue!85!black, fill=white, rounded corners=3pt, minimum width=2.3cm, minimum height=0.95cm, align=center] (bank) at (2.2,0.92) {\textbf{Savanna US Bank}\\Settled USD\\(Net of Reserve)};
  \node[draw=quotientorange!90!black, fill=white, rounded corners=3pt, minimum width=2.3cm, minimum height=0.95cm, align=center] (otc) at (5.35,0.92) {\textbf{US OTC Desk}\\USD $\to$ USDT\\($14\text{ bps}$ On-Ramp)};
  \node[draw=quotientorange!90!black, thick, fill=white, rounded corners=3pt, minimum width=2.3cm, minimum height=0.95cm, align=center] (chain) at (8.5,0.92) {\textbf{Blockchain Rail}\\USDT Transfer\\($1\text{ bp}$ Gas)};
  \node[draw=structurepurple!90!black, thick, fill=white, rounded corners=3pt, minimum width=2.05cm, minimum height=0.95cm, align=center] (hcusdt) at (11.2,0.92) {\textbf{HoneyCoin}\\USDT Wallet $\to$\\KES Pool Refill};

  \draw[->, thick, sourceblue!80!black] (pspa.south) .. controls (4.1,2.35) and (2.4,2.35) .. node[pos=0.45, fill=white, inner sep=1pt, font=\tiny] {\texttt{T+2} Wire} (bank.north);
  \draw[->, thick, quotientorange!85!black] (bank) -- (otc);
  \draw[->, thick, quotientorange!85!black] (otc) -- (chain);
  \draw[->, thick, quotientorange!85!black] (chain) -- (hcusdt);
  \draw[->, thick, structurepurple] (hcusdt.north) .. controls (11.1,2.35) and (8.6,2.35) .. node[pos=0.5, fill=white, inner sep=1pt, font=\tiny\bfseries] {Restores KES Float} (hcapi.south);
\end{tikzpicture}
\caption{Savanna Pay's four-ledger cross-border architecture. Top lane: Collection PSP captures dollars in New York while pre-funded HoneyCoin float pays M-Pesa shillings in ten seconds. Bottom lane: settled US bank cash converts via an OTC desk to \texttt{USDT} and refills HoneyCoin in thirty minutes.}
\label{fig:four-ledger-settlement}
\end{figure}

\section{Settling One 200 USD Transfer Across All Four Ledgers}

Before looking at a busy day's batch, trace \textbf{one single transfer} from end to end so every debit and credit makes complete physical sense. Recall our sign rules from Chapter~1: Assets and Expenses increase on the \textbf{Debit} (left) side, while Liabilities and Revenue increase on the \textbf{Credit} (right) side.

Grace's cousin Marcus lives in New York. On Monday morning, Marcus sends $200.00\text{ USD}$ on Savanna Pay to Mary's M-Pesa wallet in Nairobi. Here are the exact terms on Monday morning:
\begin{itemize}
  \item \textbf{Wholesale mid-market rate.} In the wholesale market, $1\text{ USD} = 130.00\text{ KES}$.
  \item \textbf{Marcus's locked customer quote.} Savanna Pay quotes $1\text{ USD} = 127.66\text{ KES}$ with $0\text{ USD}$ upfront fee. That rate is $1.80\%$ ($180\text{ bps}$, or $3.60\text{ USD}$) below mid-market. Mary is promised $200 \times 127.66 = 25{,}532.00\text{ KES}$, whose wholesale value at $130.00$ is $25{,}532 / 130.00 = 196.40\text{ USD}$.
  \item \textbf{Collection PSP terms.} Blended bank and card fee of $0.80\%$ ($80\text{ bps}$, or $1.60\text{ USD}$). Cash settles on \texttt{T+2}, with $5\%$ ($10.00\text{ USD}$) held in a 90-day rolling reserve.
  \item \textbf{OTC On-Ramp and Gas terms.} Conversion costs $14\text{ bps}$ ($0.28\text{ USD}$) at the OTC desk plus $0.02\text{ USD}$ ($1\text{ bp}$) in blockchain gas, totaling $0.30\text{ USD}$ ($15\text{ bps}$).
  \item \textbf{HoneyCoin Payout terms.} HoneyCoin converts $1\text{ USDT} = 130.00\text{ KES}$ into Kenyan float and charges a $40\text{ bps}$ ($0.80\text{ USD} = 104.00\text{ KES}$) M-Pesa payout fee.
\end{itemize}

Watch all four ledgers move in sequence.

\paragraph{Ledger Step 1 ($t = 2\text{ seconds}$): Collection PSP Captures Marcus's $200.00\text{ USD}$}
The US Collection PSP charges Marcus $200.00\text{ USD}$ and keeps its $1.60\text{ USD}$ fee. Savanna Pay does not have cash in its bank yet. Instead, the PSP owes Savanna Pay $198.40\text{ USD}$ (an Asset we debit), Savanna Pay incurs $1.60\text{ USD}$ of fee cost (an Expense we debit), and Savanna Pay now owes Marcus a $200.00\text{ USD}$ remittance delivery (a Liability we credit):

\begin{center}
\small
\begin{tabular}{llrr}
\toprule
\textbf{Step 1 Entry (Single Transfer)} & \textbf{Type} & \textbf{Debit (USD)} & \textbf{Credit (USD)} \\
\midrule
PSP Clearing Receivable & Asset ($\uparrow$) & \$198.40 & \\
Collection PSP Fee Expense & Expense ($\uparrow$) & \$1.60 & \\
Customer Remittance Obligation & Liability ($\uparrow$) & & \$200.00 \\
\bottomrule
\end{tabular}
\end{center}

\paragraph{Ledger Step 2 ($t = 8\text{ seconds}$): HoneyCoin Pays Mary's M-Pesa Instantly}
Savanna Pay calls HoneyCoin's payout API. HoneyCoin sends $25{,}532.00\text{ KES}$ (worth $196.40\text{ USD}$) to Mary's M-Pesa phone and charges its $104.00\text{ KES}$ ($0.80\text{ USD}$) rail fee. In total, Savanna Pay's pre-funded \texttt{HoneyCoin KES Float} drops by $25{,}636.00\text{ KES}$ ($197.20\text{ USD}$ exact at $130.00\text{ KES/USD}$, credited to reduce the Asset).

Why does Revenue appear right here? Because delivering Mary's shillings extinguishes our full $200.00\text{ USD}$ obligation to Marcus (debited to zero) using only $196.40\text{ USD}$ worth of shillings. That $3.60\text{ USD}$ gap is our earned FX spread (credited to Revenue), alongside the $0.80\text{ USD}$ HoneyCoin fee (debited to Expense):

\begin{center}
\small
\begin{tabular}{llrr}
\toprule
\textbf{Step 2 Entry (Single Transfer)} & \textbf{Type} & \textbf{Debit (USD)} & \textbf{Credit (USD)} \\
\midrule
Customer Remittance Obligation & Liability ($\downarrow$) & \$200.00 & \\
HoneyCoin Payout Rail Expense & Expense ($\uparrow$) & \$0.80 & \\
HoneyCoin KES Float ($25{,}636\text{ KES}$) & Asset ($\downarrow$) & & \$197.20 \\
Gross FX Spread Revenue & Revenue ($\uparrow$) & & \$3.60 \\
\bottomrule
\end{tabular}
\end{center}

Check both columns: Debits equal $200.80\text{ USD}$, and Credits equal $197.20 + 3.60 = 200.80\text{ USD}$. Eight seconds after Marcus tapped Send, Mary has her shillings and Marcus's obligation sits at zero.

\paragraph{Ledger Step 3 (\texttt{T+2 Morning}): Collection PSP Wires Cash (Minus Rolling Reserve)}
Two business days later, the Collection PSP settles the $198.40\text{ USD}$ it owed us. It holds $10.00\text{ USD}$ ($5\%$ of $200\text{ USD}$) inside a 90-day rolling reserve Asset and wires the remaining $188.40\text{ USD}$ into Savanna Pay's US bank account, clearing \texttt{PSP Clearing Receivable} down to zero:

\begin{center}
\small
\begin{tabular}{llrr}
\toprule
\textbf{Step 3 Entry (\texttt{T+2} Bank Clear)} & \textbf{Type} & \textbf{Debit (USD)} & \textbf{Credit (USD)} \\
\midrule
US Settlement Bank (USD) & Asset ($\uparrow$) & \$188.40 & \\
PSP Rolling Reserve Receivable & Asset ($\uparrow$) & \$10.00 & \\
PSP Clearing Receivable & Asset ($\downarrow$) & & \$198.40 \\
\bottomrule
\end{tabular}
\end{center}

\begin{recoverybox}
\textbf{Pause and think:} Look closely at that $10.00\text{ USD}$ rolling reserve in Step~3. Did the rolling reserve reduce Savanna Pay's profit? No: it moved $10.00\text{ USD}$ from one Asset (\texttt{PSP Clearing Receivable}) to another Asset (\texttt{PSP Rolling Reserve Receivable}), which releases in $90$ days. Yet look at the cash flow: Savanna Pay received $188.40\text{ USD}$ in bank cash today while it spent $197.20\text{ USD}$ of Nairobi float on Monday. Fast-growing card corridors need extra working capital to bridge that $90\text{-day}$ rolling reserve gap.
\end{recoverybox}

\paragraph{Ledger Step 4 (\texttt{T+2 + 30 Minutes}): OTC USDT On-Ramp and HoneyCoin Float Refill}
To put back the exact $25{,}636.00\text{ KES}$ ($197.20\text{ USD}$) spent from Nairobi float, Savanna Pay wires $197.48\text{ USD}$ from its US bank to its US OTC desk. The OTC desk deducts $0.28\text{ USD}$ fee and sends $197.20\text{ USDT}$, which Savanna Pay transfers on-chain to HoneyCoin (burning $0.02\text{ USD}$ from \texttt{Crypto Gas Inventory}) and swaps into $25{,}636.00\text{ KES}$:

\begin{center}
\small
\begin{tabular}{llrr}
\toprule
\textbf{Step 4 Entries (30-Min Refill)} & \textbf{Type} & \textbf{Debit (USD)} & \textbf{Credit (USD)} \\
\midrule
HoneyCoin KES Float ($25{,}636\text{ KES}$) & Asset ($\uparrow$) & \$197.20 & \\
OTC On-Ramp \& Gas Expense & Expense ($\uparrow$) & \$0.30 & \\
US Settlement Bank (USD) & Asset ($\downarrow$) & & \$197.48 \\
Crypto Gas Inventory (TRX/ETH) & Asset ($\downarrow$) & & \$0.02 \\
\bottomrule
\end{tabular}
\end{center}

Add up Revenue minus Expenses on Marcus's single $200.00\text{ USD}$ transfer:
\[
  +3.60\text{ (FX spread)} - 1.60\text{ (PSP)} - 0.80\text{ (HoneyCoin)} - 0.30\text{ (OTC + gas)} = +0.90\text{ USD } (+45\text{ bps}).
\]
Now add up Savanna Pay's net Asset changes across all four steps: $+188.40 - 197.48 = -9.08\text{ USD}$ in US Bank, $+10.00\text{ USD}$ in PSP Rolling Reserve, $0.00\text{ USD}$ net change in HoneyCoin KES float (fully refilled), and $-0.02\text{ USD}$ in Gas Inventory. Their sum is $-9.08 + 10.00 - 0.02 = +0.90\text{ USD}$. Income statement and balance sheet agree to the penny.

\section{Scaling to a 50,000 USD Daily Batch and Breakeven Math}

In real operations, Nia does not send tiny $197\text{ USD}$ wires one by one. Savanna Pay lets individual payouts fire every few seconds out of HoneyCoin's pre-funded float, and then settles the wholesale bank-to-\texttt{USDT}-to-HoneyCoin leg once or twice a day in batches.

Let us scale Marcus's transfer by $250\times$ to trace a daily batch of $250$ US-to-Kenya transfers totaling $50{,}000.00\text{ USD}$.

\paragraph{Scenario Parameters}
\begin{itemize}
  \item \textbf{Batch Volume.} Across Monday, $250$ US senders pay $200.00\text{ USD}$ each, giving $\text{GTV} = 50{,}000.00\text{ USD}$ at $127.66\text{ KES/USD}$. Relative to the $130.00$ mid-market rate, Savanna Pay delivers $6{,}383{,}000\text{ KES}$ (worth $49{,}100.00\text{ USD}$) and earns $900.00\text{ USD}$ ($180\text{ bps}$) in gross FX spread.
  \item \textbf{Collection PSP Cost ($80\text{ bps}$).} The PSP deducts $400.00$ \texttt{USD} in fees and settles $49{,}600.00$ \texttt{USD} net value into Savanna Pay's US bank account.
  \item \textbf{USDT Bridge and HoneyCoin Payout Cost ($55\text{ bps}$).} On-ramping costs $70.00\text{ USD}$ ($14\text{ bps}$), blockchain gas costs $5.00\text{ USD}$ ($1\text{ bp}$), and HoneyCoin payout fees cost $200.00\text{ USD}$ ($40\text{ bps}$, or $26{,}000\text{ KES}$). Combined bridge and payout costs total $275.00\text{ USD}$.
  \item \textbf{Net Batch Profit ($+45\text{ bps}$).} Subtracting $400.00$ and $275.00$ from $900.00\text{ USD}$ leaves $+225.00\text{ USD}$ net profit after restoring all $49{,}300.00\text{ USD}$ ($6{,}409{,}000\text{ KES}$) spent from HoneyCoin float (Figure~\ref{fig:stage-balance-timeline}).
\end{itemize}

\begin{figure}[htbp]
\centering
\begin{tikzpicture}[x=1cm,y=1cm,>=Latex,font=\scriptsize]
  \draw[thick, ->] (2.8,4.15) -- (12.2,4.15);
  \node[font=\bfseries\scriptsize] at (4.0,4.42) {Stage 1 ($t=2\text{s}$)};
  \node[font=\bfseries\scriptsize] at (6.4,4.42) {Stage 2 ($t=10\text{s}$)};
  \node[font=\bfseries\scriptsize] at (8.9,4.42) {Stage 3 (Bank)};
  \node[font=\bfseries\scriptsize] at (11.3,4.42) {Stage 4 (+30m)};

  \node[anchor=east, font=\bfseries\scriptsize] at (2.65,3.45) {PSP Receivable};
  \node[anchor=east, font=\bfseries\scriptsize] at (2.65,2.60) {US Bank (USD)};
  \node[anchor=east, font=\bfseries\scriptsize] at (2.65,1.75) {USDT Transit};
  \node[anchor=east, font=\bfseries\scriptsize] at (2.65,0.90) {HoneyCoin KES};
  \node[anchor=east, font=\bfseries\scriptsize, text=targetgreen!80!black] at (2.65,0.12) {Net Profit ($\Delta E$)};

  \foreach \y in {0.50,1.32,2.17,3.02,3.88} {
    \draw[gray!25] (2.8,\y) -- (12.2,\y);
  }

  \node[fill=sourceblue!12, rounded corners=2pt, inner sep=2.5pt] at (4.0,3.45) {$+$\$49,600};
  \node[text=gray] at (6.4,3.45) {$+$\$49,600};
  \node[fill=gray!12, rounded corners=2pt, inner sep=2.5pt] at (8.9,3.45) {\$0};
  \node[text=gray] at (11.3,3.45) {\$0};

  \node[text=gray] at (4.0,2.60) {\$0};
  \node[text=gray] at (6.4,2.60) {\$0};
  \node[fill=sourceblue!12, rounded corners=2pt, inner sep=2.5pt] at (8.9,2.60) {$+$\$49,600};
  \node[fill=sourceblue!12, rounded corners=2pt, inner sep=2.5pt] at (11.3,2.60) {$+$\$230};

  \node[text=gray] at (4.0,1.75) {$0\text{ USDT}$};
  \node[text=gray] at (6.4,1.75) {$0\text{ USDT}$};
  \node[fill=quotientorange!15, rounded corners=2pt, inner sep=2.5pt] at (8.9,1.75) {(Swap $49{,}300$)};
  \node[fill=gray!12, rounded corners=2pt, inner sep=2.5pt] at (11.3,1.75) {$0\text{ USDT}$};

  \node[text=gray] at (4.0,0.90) {\$0};
  \node[fill=trapred!12, rounded corners=2pt, inner sep=2.5pt] at (6.4,0.90) {$-$\$49,300};
  \node[text=gray] at (8.9,0.90) {$-$\$49,300};
  \node[fill=structurepurple!12, rounded corners=2pt, inner sep=2.5pt] at (11.3,0.90) {\$0 (Refilled)};

  \node[text=trapred!85!black] at (4.0,0.12) {$-$\$400};
  \node[fill=targetgreen!14, rounded corners=2pt, inner sep=2.5pt, font=\bfseries\scriptsize] at (6.4,0.12) {$+$\$300};
  \node[text=targetgreen!85!black, font=\bfseries\scriptsize] at (8.9,0.12) {$+$\$300};
  \node[fill=targetgreen!18, rounded corners=2pt, inner sep=2.5pt, font=\bfseries\scriptsize, text=targetgreen!75!black] at (11.3,0.12) {$+$\$225};
\end{tikzpicture}
\caption{Cumulative balance-sheet changes across the $50{,}000\text{ USD}$ batch. Wiring $49{,}370\text{ USD}$ from US Bank swaps into $49{,}300\text{ USDT}$ and refills HoneyCoin KES float completely to zero delta, leaving $+230\text{ USD}$ in US Bank minus $5\text{ USD}$ gas $= +225\text{ USD}$ net profit.}
\label{fig:stage-balance-timeline}
\end{figure}

\subsection{The Corridor Breakeven Formula (Including Chargebacks)}

Every corridor Savanna Pay launches must pass one clear unit-economics test. Let $s_{\text{fx}}$ be the customer FX spread (plus any explicit fee) in basis points. Subtract the four basis-point cost layers: $c_{\text{collect}}$ (Collection PSP), $c_{\text{onramp}}$ (OTC on-ramp plus gas), $c_{\text{honeycoin}}$ (local payout rail), and $c_{\text{chargeback}}$ (expected card chargeback losses). Earning positive margin requires:
\[
  s_{\text{fx}} > c_{\text{collect}} + c_{\text{onramp}} + c_{\text{honeycoin}} + c_{\text{chargeback}}.
\]

\begin{trapbox}
\textbf{Reversible Card Inflows vs Irreversible Mobile-Money Payouts.} Never ignore $c_{\text{chargeback}}$ on debit or credit card corridors. Once Mary receives $25{,}532\text{ KES}$ on M-Pesa and cashes out at Joseph's kiosk, that payout can never be pulled back. Meanwhile, a thief in New York using a stolen debit card gives the real cardholder up to $120\text{ days}$ to file a bank chargeback. One disputed $200\text{ USD}$ card transfer claws back all $200\text{ USD}$ of principal plus a $15\text{ USD}$ penalty ($215\text{ USD}$ lost), erasing the $+0.90\text{ USD}$ profit from $239$ honest transfers. When card fees cost $145\text{ bps}$ and fraud costs $25\text{ bps}$, your breakeven spread starts at $145 + 15 + 40 + 25 = 225\text{ bps}$ before earning a cent.
\end{trapbox}

\begin{intuitionbox}
What if a Kenyan electronics importer wants to pay a supplier in New York or Shenzhen? The exact same HoneyCoin pipeline runs in reverse. HoneyCoin acts as the \textbf{Collection Partner} in Nairobi, collecting shillings via M-Pesa STK push or a local virtual bank account into Savanna Pay's Kenyan float. Savanna Pay swaps those shillings on HoneyCoin's OTC desk into \texttt{USDT} and off-ramps to \texttt{USD} for the supplier's bank. Better still, Savanna Pay can net those shillings directly against incoming US-to-Kenya remittances with zero blockchain or OTC fees.
\end{intuitionbox}

\begin{walkbox}
\textbf{Quick pricing check:} Savanna Pay opens a pure debit-card corridor to Uganda where card collection costs $c_{\text{collect}} = 150\text{ bps}$, on-ramp plus gas costs $c_{\text{onramp}} = 15\text{ bps}$, HoneyCoin MTN MoMo payout costs $c_{\text{honeycoin}} = 45\text{ bps}$, and expected card chargebacks cost $c_{\text{chargeback}} = 20\text{ bps}$. What customer FX spread $s_{\text{fx}}$ is needed to earn a $+50\text{ bps}$ net margin? Add the four cost layers ($150 + 15 + 45 + 20 = 230\text{ bps}$ breakeven) and add $50\text{ bps}$ target margin to get $s_{\text{fx}} = 280\text{ bps}$ ($2.80\%$).
\end{walkbox}

Savanna Pay now has money moving across four independent systems: a US Collection PSP, a US bank, a public blockchain, and HoneyCoin's mobile-money pools. At 08:00 AM tomorrow morning, how do Nia, Sam, and Tunde prove -- down to the exact cent and shilling -- that zero money leaked overnight, and what reports keep their partner banks happy? That is the story of Chapter~\ref{ch:reconciliation-and-reporting}.

\section*{Exercises}
\addcontentsline{toc}{section}{Exercises}

\begin{exercise}[Tracing a Single 100 USD Card-to-MoMo Remittance]
A sender in Chicago pays $100.00\text{ USD}$ by debit card to send money to a relative's MTN MoMo wallet in Kampala. Wholesale mid-market is $\texttt{USD/UGX} = 3{,}700.00$, and Savanna Pay quotes $3{,}607.50\text{ UGX/USD}$ with $0\text{ USD}$ upfront fee. The Collection PSP charges $1.50\% + 0.30\text{ USD}$ and holds a $5\%$ rolling reserve at \texttt{T+2}. HoneyCoin converts $1\text{ USDT} = 3{,}700.00\text{ UGX}$ and charges a $1{,}850\text{ UGX}$ payout fee. OTC on-ramp plus blockchain gas costs $0.15\text{ USD}$ total.
\begin{enumerate}
  \item Write the four double-entry journal batches for this single $100.00\text{ USD}$ transfer across capture, instant MTN MoMo payout, \texttt{T+2} bank settlement, and 30-minute \texttt{USDT} float refill.
  \item Compute Savanna Pay's exact net profit or loss on this transfer in both dollars and basis points. Is a $250\text{ bps}$ FX spread wide enough for a $100\text{ USD}$ debit-card ticket?
\end{enumerate}
\end{exercise}

\begin{exercise}[Bitcoin Volatility vs USDT Stability in Transit]
Savanna Pay needs to move a $100{,}000\text{ USD}$ wholesale batch to refill its HoneyCoin payout pool on a corridor earning $50\text{ bps}$ ($500\text{ USD}$) before bridge costs.
\begin{enumerate}
  \item \textbf{Option A (Unhedged Base-Layer BTC):} Buy \texttt{BTC} at $10\text{ bps}$ fee, wait $45\text{ minutes}$ for confirmations ($4.00\text{ USD}$ miner fee), and sell \texttt{BTC} into local float at $10\text{ bps}$ fee. Compute Savanna Pay's net batch profit if \texttt{BTC} stays flat, and its net loss if \texttt{BTC} drops by $1.2\%$ during those $45\text{ minutes}$.
  \item \textbf{Option B (USDT Bridge):} On-ramp \texttt{USD} to \texttt{USDT} at $12\text{ bps}$ fee, pay $1.50\text{ USD}$ in native \texttt{TRX} gas, confirm in $30\text{ seconds}$, and swap into local float at $0\text{ bps}$ par. Compute the exact net profit in dollars and basis points.
\end{enumerate}
\end{exercise}

\begin{exercise}[Chargeback Breakeven Sensitivity]
On Savanna Pay's $50{,}000\text{ USD}$ daily batch ($250$ transfers of $200\text{ USD}$), suppose $s_{\text{fx}} = 180\text{ bps}$, $c_{\text{collect}} = 80\text{ bps}$, $c_{\text{onramp}} = 15\text{ bps}$, and $c_{\text{honeycoin}} = 40\text{ bps}$, giving $+225.00\text{ USD}$ ($+45\text{ bps}$) before fraud.
\begin{enumerate}
  \item If each disputed $200\text{ USD}$ card transfer loses the full $200\text{ USD}$ payout plus a $15\text{ USD}$ PSP chargeback fee ($215\text{ USD}$ total loss), how many fraudulent transfers out of $250$ completely wipe out the day's $225.00\text{ USD}$ profit?
  \item Express that critical fraud rate as a percentage of batch volume and in basis points $c_{\text{chargeback}}^*$.
\end{enumerate}
\end{exercise}

\begin{exercise}[Reverse Corridor: Collecting KES via HoneyCoin to Pay a US Supplier]
A Nairobi solar-panel distributor needs Savanna Pay Business to deliver $20{,}000.00\text{ USD}$ net via domestic US wire to a manufacturer in Texas. HoneyCoin charges $0.60\%$ to collect \texttt{KES} locally and quotes an outbound OTC ask rate of $130.65\text{ KES/USDT}$ when wholesale mid-market is $130.00\text{ KES/USD}$. On-chain gas costs $2.00\text{ USD}$, and Savanna Pay's US OTC off-ramp charges $10\text{ bps}$ plus a $15.00\text{ USD}$ wire fee. Compute the exact \texttt{KES} invoice total Savanna Pay must collect from the Nairobi distributor to earn a $+60\text{ bps}$ ($120.00\text{ USD}$) net profit.
\end{exercise}
